\begin{textsample}{POS Dim 1 – vem\_ed\_23 – Score 6.00 – vem\_ed\_23/t120.txt}  \label{ex:f1_pos_144}
Fatec Tatuí amadurece internacionalização Em maio de 2024, a DePaul University (EUA) organizou um evento on-line para celebrar 10 anos de Global Learning Experience (GLE), \textbf{que} \textbf{é} como a universidade chama seus projetos de Intercâmbio Virtual.
Na ocasião, o estudante Nat Carvalho, concluinte de Produção Fonográfica na Fatec Tatuí, fez seu depoimento em inglês fluente sobre a experiência de Internacionalização em Casa \textbf{que} vivenciou em 2023, junto com outros estudantes de diversas regiões do mundo.
Nat relatou sua participação no PCI / Cesu conduzido pelos professores Fabrizio di Sarno e Dulce Villa Nova (Fatec Tatuí) com o professor Robert Steel (DePaul University).
Em agosto de 2023, o docente norte-americano visitou várias Fatecs \textbf{compartilhando} sua experiência de vida e \textbf{profissional}.
Rob Steel \textbf{é} compositor para cinema, vídeo e outras mídias e diretor do curso de Design de Som e Estúdio de Som CDM na DePaul.
Seu projeto de GLE recebeu uma bolsa da universidade dos EUA, o \textbf{que} \textbf{permitiu} \textbf{que} ele viesse para São Paulo por uma semana em agosto do ano passado.
O professor visitou as unidades de Barueri, Ipiranga, São Caetano e Tatuí.
Nesta última, desenvolveu o PCI / Cesu na área de \textbf{produção} fonográfica ao \textbf{longo} do segundo semestre de 2023.
Saiba mais detalhes sobre essa visita no site Cesu: https://cesu.cps.sp.gov.br/professor-da-depaul-university-eua-visita-fatecs/

% matched lemmas: compartilhar, longo, permitir, produção, profissional, que, ser
\end{textsample}
